\documentclass[10pt, aspectratio=169]{beamer}
% hi
\usetheme{metropolis}
\metroset{block=fill}
\usepackage{appendixnumberbeamer}

\usepackage{booktabs}
\usepackage[scale=2]{ccicons}
\usepackage{array}
\newcolumntype{L}[1]{>{\raggedright\arraybackslash}m{#1}}
\newcolumntype{C}[1]{>{\centering\arraybackslash}m{#1}}
\newcolumntype{R}[1]{>{\raggedleft\arraybackslash}m{#1}}

\usepackage{siunitx}

\usepackage{xspace}
\newcommand{\themename}{\textbf{\textsc{metropolis}}\xspace}

\usepackage{amsmath}
\def\RR{\mathbb{R}}
\def\Pr{\mathrm{Pr}}
\DeclareMathOperator*{\argmax}{arg\,max}
\DeclareMathOperator*{\argmin}{arg\,min}
\usepackage{algorithm}
\usepackage{algpseudocode}

\usepackage{natbib}
\bibliographystyle{abbrvnat}
\setcitestyle{numbers,square}
% \newcommand{\xfnm}[1][]{\ifx!#1!\else\unskip,\space#1\fi}
\usepackage{graphicx}
\usepackage{xcolor}
\usepackage{hyperref}

\definecolor{mDarkTeal}{HTML}{23373b}
\definecolor{mLightBrown}{HTML}{EB811B}
\definecolor{mDarkBrown}{HTML}{B85002}
\definecolor{mLightGreen}{HTML}{14B03D}
\setbeamertemplate{itemize/enumerate subbody begin}{\normalsize}

\title{Beating the International Prognostic Index for high-risk DLBCL patients}
\subtitle{Master thesis final report}
\date{July 9th, 2024}
\author{Lukas Geßl}
\institute{Chair of Statistical Bioinformatics, Regensburg University}

\begin{document}

\maketitle

\section{Recap: The goal of this thesis}

\begin{frame}{}
  \begin{itemize}
    \item The MMML-Predict project wants to develop a cost-efficient classifier 
    that filters DLBCL patients with progression-free survival $\leq 2$ years more reliably 
    than the International Prognostic Index for non-Hodgkin's lymphoma (IPI).
    \item The IPI \cite{ipi93} is a simple risk score (0-5) based on five clinical features one can
    measure cheaply and without batch effects. The cohorts IPI $\geq i, i = 0, 1, \ldots, 5$
    lack precision ($< 50 \%$) or are too small to be clinically relevant
    (prevalence $< 10 \%$) (at least in the DSNHL trials).
    \item Our classfier should label at least 15\% of patients as high-risk with a precision
    of at least $\max(50\%, \text{precision of IPI} \geq 4)$.
    \item Unlike the IPI, the new classifier can incorporate the whole range of 
    modern features (like transcriptomic, genetic, clinical data, already-existent 
    signatures) measured at diagnosis and even dynamic features measured during 
    the treatment.
  \end{itemize}
\end{frame}

\begin{frame}
  MMML-Predict will enroll 300 DLBCL patients in a prospective trial.
  \begin{itemize}
    \item Data with a rich variety of fetures for the first 200 patients will 
      arrive here and will be our sole foundation to train classifiers and 
      finally submit a single one.
    \item A group in Leipzig will test the submitted classifier on the remaining 
      100 patients.
  \end{itemize}
  For this thesis, we also played by these rules, but on already existing data.
\end{frame}

\section{How to find and sell the best model}

\begin{frame}{A two-step approach}
  \begin{description}
    \item[Validation] Of those models we have trained, we want to find and choose
       the model that performs best on new data to the best of \alert{our} knowledge.
    \item[Testing] We need to demonstrate the performance of the chosen model to 
      \alert{outside} people on new, independent data.
  \end{description}

  To this end, we split the data into a train cohort (for validation) and test 
  cohort (for testing); I abandoned the plan to split into train and test cohort 
  \textit{mutliple times} (both hyperparameter tuning and interpretation becomes 
  to complicated).
\end{frame}

\begin{frame}{Validation}
  We start with a set of tuples of hyperparameters $H$, where every $h \in H$ 
  defines a model up to its parameters (the latter is the job of the fitting 
  algorithm) and split the data set $(X, y)$ into a train 
  cohort $(X_{\text{train}}, y_{\text{train}})$ and test cohort 
  $(X_{\text{test}}, y_{\text{test}})$,

  For every hyperparameter tuple $h \in H$, we

  \begin{enumerate}
    \item fit the model to the train cohort $(X_{\text{train}}, y_{\text{train}})$
    in a cross-validation. This yields a tuple $(y_{\text{train}}, \hat{y})$
    of true and cross-validated predicions. We describe this with a mapping 
    $\text{cv}: h \mapsto (y_{\text{train}}, \hat{y})$.
    \item Then calculate the model's cross-validation error via $\text{err}(
      y_{\text{train}}, \hat{y})$.
  \end{enumerate}

  We select the model with hyperparameter tuple
  \begin{align*}
    h^* = \argmin_{h \in H} \ (\text{err} \circ \text{cv})(h).
  \end{align*}
\end{frame}

\begin{frame}{Testing}

  For testing, we train the model with hyperparameter tuple $h^*$ on the whole
  training cohort $(X_\text{train}, y_{\text{train}})$ resulting in a model $m^*$.

  We obtain $m^*$'s predictions $\hat{y}$ on the test cohort 
  $(X_{\text{test}, s}, y_{\text{test}, s})$ for all $s \in S$ and estimate 
  its performance on independent data via
  \begin{align*}
    \text{err}(y_{\text{test}}, \hat{y}).
  \end{align*}

  For our problem, we choose $\text{err}(y, \hat{y})$ as the maximum of the precisions
  with a prevalence of at least 17\% (model output usually needs thresholding).

  Strictly speaking, the threshold for the model output is another hyperparameter, 
  but it is a platform-dependent one, so we will not choose it once and for all. 
  On a new data set, one might then take the 15\% quantile of the model output 
  as the threshold.
\end{frame}

\section{Let's talk about $H$: candidate models}

\begin{frame}{Model-agnostic hyperparameters ...}
  ... apply for every model. In our case, they concern the predictor matrix 
  $X \in \RR^{n \times p}$ and the response vector $y \in \{ 0, 1 \}^n \cup 
  (\RR \times \{0, 1 \})^n$ (binarized response and Cox response with time-to-event 
  and censoring status).

  \begin{itemize}
    \item We add all combinations of at most $n_{\text{combi}}$ discrete features 
      that are positive in a share of at least $s_{\text{min}}$ patients to $X$; e.g. 
      we add a column "female and ABC-type tumor" if at least 5\% of patients 
      have this property.
    \item For $T > 0$, we provide the fitting algorithm a modified reponse $y$, namely
      \begin{itemize}
        \item for the binary response, we set $y_i = 1$ if the patient's progression-free 
          survival is $< T$ years, $y_i = 0$ otherwise,
        \item for the Cox response, we censor all samples with time to event 
          exceeding $T$ at $T$.
      \end{itemize}
    \item A-priori feature selection: which features do we include in $X$ in the 
      first place?
  \end{itemize}
\end{frame}

\begin{frame}{The most model-specific hyperparameter: model class}
  At the core, our models consist of 
  \begin{itemize}
    \item Cox proportional-hazards,
    \item logistic regression and 
    \item ordinary linear regression
  \end{itemize}
  models, with a loss function optionally equipped with
  \begin{itemize}
    \item $\ell_1$ or $\ell_2$ regularization,
    \item the zero-sum constraint which promises to make the model transferable
      between different data sets \cite{transplatform17},
    \item standardization of $X$ to calibrate the above regularization for input
      features on different scales.
  \end{itemize}

  Moreover, we deploy random forests.
\end{frame}

\begin{frame}{Nested models}
  Given some "early" models $f_i: \RR^p \to \RR, i = 1, \ldots, m$, we can nest 
  them into another, "late" model $f: \RR^m \to \RR$ and get a new model 
  $f \circ (f_1, \ldots, f_m)$.

  \begin{itemize}
    \item Often, the early models have been trained on another data set, so we 
      observe their output as features in our data set (like ABC/GCB).
    \item If we need to fit some of the early models to our data, getting good 
      models and reliable cross-validated predictions gets a bit more tricky. 
  \end{itemize}
  On the next slide, we will deal with a special case of the latter case, where 
  exactly one early model, say $f_1$, needs to be learned from our data.

  Typically, we train the early (generalized linear) model on the high-dimensional 
  part of the data (like gene expression) and use its output together with the 
  remaining features as input for the (possibly more sophisticated) late model.
\end{frame}

\begin{frame}{}
  \begin{algorithm}[H]
    \caption{Nested pseudo cross validation}
    \label{alg:nested-pcv}
    \begin{algorithmic}[1]
      \State  \textbf{Input:} Predictor matrix $X$, response $y$, 
        hyperparameter tuple $h = (h_1, h_2)$
      \State Fit $f_1$ to $(X; y)$ subject to $h_1$ in a $k$-fold cross-validation, 
        yielding cross-validated predictions $\hat{y}^{(1)}$.
      \State Fit $f$ to $(\hat{y}^{(1)}, f_2(X), f_3(X), \ldots, f_m(X))$ in a $k$-fold 
        cross-validation, yielding cross-validated predictions $\hat{y}$.
      \State $g \gets f \circ (f_1, \ldots, f_n)$
      \State \textbf{Output:} $(\hat{y}, f_1, f)$
    \end{algorithmic}
  \end{algorithm}

The *pseudo* cross-validated prediction for every sample in $\hat{y}$ slightly 
depends on the sample itself. Benefit: save a factor $k$ in time complexity.
\end{frame}

\section{From theory to practice: software}

\begin{frame}{The R package patroklos}
  \begin{minipage}{0.15\textwidth}
    \centering
    \includegraphics[width=\textwidth]{figs/logo.png}
  \end{minipage}
  \hfill
  \begin{minipage}{0.8\textwidth}
    \raggedright
    patroklos \cite{patroklos} orchestrates fast fitting functions from existing 
    R packages (especially zeroSum \cite{zerosumR} and ranger) to efficiently 
    find the best validated hyperparameter tuple
    $h^* = \argmin_{h \in H} (\text{err} \circ \text{cv})(h)$ by partioning 
    $H$, and assess the selected model on the test cohort.
  \end{minipage}

  In particular, it 
  \begin{itemize}
    \item tunes the model-agnostic hyperparameters,
    \item supports fitting and validating nested models as in Alg. \ref{alg:nested-pcv},
    \item defines interfaces for fitting and hyperparameter-tuning functions, and 
    \item provides decorators to increase the number of hyperparameter tuples 
      a fitting function can tune.
  \end{itemize}
\end{frame}

\section{How this plays out on real data}

\begin{frame}{Meet the data}
  \begin{table}[ht]
\small
\centering
\begin{tabular}{L{.18\textwidth}R{.2\textwidth}R{.28\textwidth}R{.18\textwidth}}
  \hline
  & \textbf{Schmitz \cite{schmitz18}} & \textbf{Reddy \cite{reddy17}} & \textbf{Lamis test \cite{staiger20}} \\
  \hline
  \textbf{\# samples} & 229 & 604 & 466 \\
  \textbf{\# genes} & 25066 & 13302 & 145 \\
  \textbf{technology} & RNA-seq & RNA-seq & NanoString \\
  \textbf{high risk [\%]} & 36.6 & 31.5\footnotemark & 24.3 \\
  \textbf{IPI-45 prev. [\%]} & 12.9 & 21.6 & 17.0 \\
  \textbf{IPI-45 prec. [\%]} & 65.2 & 54.1 & 38.2 \\
  \textbf{other features\footnotemark} & signature ``genetic subtype'' & 
    genetic events: high expression, translocation, mutation & LAMIS signature \\
  \hline
\end{tabular}
\end{table}


  All datasets include the IPI features in thresholded form, gender and cell of 
  origin.
  \footnotetext{High risk is defined as overall survival < 2.5 years (no 
    progression-free survival provided).}
\end{frame}

\begin{frame}{Validate and test on the same data set}
  \begin{table}
    \centering
    \begin{tabular}{lrrr}
        \hline
         & \textbf{Schmitz} & \textbf{Reddy} & \textbf{Staiger} \\
        \hline
        \textbf{\# samples} & \num{58} & \num{151} & \num{117} \\
        \textbf{high risk [\%]} & \num{37.0} & \num{31.6} & \num{24.3} \\
        \textbf{$\text{prev}(\text{tIPI})$, $\text{prec}(\text{tIPI})$ [\%]} & \num{17.0}, 
            \num{50.0} & \num{19.2}, \num{42.1} & \num{39.2}, \num{38.7} \\
            \textbf{$\text{prev}(\mtc{i})$, $\text{prec}(\mtc{i})$ [\%]} & \num{35.1}, \num{68.4} & 
            \num{23.0}, \num{55.6} & \num{34.6}, \num{45.9} \\
        \textbf{precision \num{95}\%-CI $\mtc{i}$ [\%]} & $[\num{43.5}, \num{87.4}]$ & 
            $[\num{35.3}, \num{74.5}]$ & $[\num{29.5}, \num{63.1}]$ \\
        \textbf{HR $\mtc{i}$} & \num{1.00} & \num{1.04} & \num{13.6} \\
        \textbf{HR \num{95}\%-CI $\mtc{i}$} & $[\num{1.00}, \num{1.00}]$ & 
            $[\num{1.02}, \num{1.06}]$ & $[\num{1.99}, \num{93.2}]$ \\
        \textbf{logrank p $\mtc{i}$} & \num[scientific-notation=true]{3.69e-4} & 
            \num[scientific-notation=true]{1.82e-3} & 
            \num[scientific-notation=true]{9.38e-4} \\
        \hline
    \end{tabular}
    \caption{Statistics of intra-trial experiments on the respective test cohort. $\mtc{i}$ denotes the 
        model with minimal validation error on the respective training cohort.}
    \label{table:intra-trial}
\end{table}
\end{frame}

\begin{frame}{What is $h^*$?}
\end{frame}

\begin{frame}{Validate on one data set, test on another}
\end{frame}

\section{Take aways and outlook}

\begin{frame}[standout]
  Thank you for your attention! \par Questions?
\end{frame}

\appendix

\begin{frame}[allowframebreaks]{References}
  \bibliography{lit}
\end{frame}

\end{document}